\documentclass[11pt]{article}
\usepackage[margin=1.2in]{geometry}
\usepackage{amsmath,amssymb,amsthm}
\usepackage{hyperref}
\sloppy

\newtheorem{theorem}{Theorem}
\newtheorem{corollary}[theorem]{Corollary}
\newtheorem{lemma}[theorem]{Lemma}
\newtheorem{conjecture}[theorem]{Conjecture}
\newtheorem{proposition}[theorem]{Proposition}
\theoremstyle{remark}
\newtheorem{remark}[theorem]{Remark}

\newcommand{\E}{\mathbb{E}}
\newcommand{\Q}{\mathbb{Q}}

\title{Roots of $d$-matching polynomials and the spectrum of the universal cover:\\
a closed form at first Betti number one, and a refuted conjecture}
\author{Vico Bonfioli\thanks{Lean~4 sources: \texttt{RamaLean/} in the \texttt{rama-notes} repository. Correspondence: \texttt{vicobonfioli@gmail.com}.}}
\date{August 2026}

\begin{document}
\maketitle

\begin{abstract}
Hall, Puder and Sawin place the roots of the $d$-matching polynomial $\mu_{d,G}$ in the interval
$[-\rho,\rho]$, $\rho$ the spectral radius of the universal cover $T$. We conjectured that they
lie in $\operatorname{spec}(T)$ itself, a proper subset of that interval whenever $T$ has a spectral
gap: no root ever enters a gap. This strengthens their theorem and is the matching-polynomial form
of their Question~6.3.

We prove it for every connected $G$ of first Betti number one and every $d$, by giving $\mu_{d,G}$
in closed form.

In general it is false. The counterexample is due to Chris Hall \cite{HallCE}: a bipartite graph on
$41$ vertices with $\mu_G=x^{21}(x^4-11x^2+25)^4(x^2-5)(x^2-11)$, where $\sqrt5$ lies in an internal
gap of $\operatorname{spec}(T)$. We give the mechanism. A separator of size $k$ met by $p$ identical
branches forces $A^{\,p-k}\mid\mu_G$, where $A$ is the branch's own matching polynomial
(Lemma~\ref{lem:divis}); Hall's $\sqrt5$ is the top matching root of the star $K_{1,5}$, and a
star's leaves have degree one in the divisor however large their degree is in $G$. Hypotheses on
$G$ therefore cannot reach the offending root. We show this is not a figure of speech by refuting
both repairs that have been proposed: minimum degree three fails on $14$ vertices, with
$\mu_G=x^4(x^2-3)(x^8-24x^6+171x^4-396x^2+270)$ and $\sqrt3$ outside $\operatorname{spec}(T)$, and
$3$-connectivity fails on $23$. Any repair must constrain the divisor rather than the graph.

Everything asserted is machine-checked in Lean~4/Mathlib unless labelled otherwise. Supplementary
material is in \cite{Supp}; the biregular case, where the conjecture survives, is \cite{PartII}.
\end{abstract}

\medskip
\noindent\textbf{2020 Mathematics Subject Classification.} Primary 05C31; Secondary 05C50,
05C70, 68V20.

\noindent\textbf{Keywords.} matching polynomial, $d$-matching polynomial, universal cover,
spectral gap, graph covering, interlacing families, Lean~4, formalization.

\noindent\textbf{Code and data availability.} All computations, the Lean~4/Mathlib sources, and
the scripts backing every numerical statement are in the \texttt{rama-notes} repository. Each
numerical claim names the script that produced it; every named script is executed and its output
compared against a recorded baseline by \texttt{code/cite\_check.py}, and every Lean declaration
cited is checked for its axiom base by \texttt{code/audit\_papers.py}.

\paragraph{Status conventions.} Claims here are of three kinds, and the text says which each time
rather than leaving it to the reader. A result described as \emph{formalized}, or cited by a Lean
declaration name, is machine-checked in Lean~4/Mathlib: the sources are in \texttt{RamaLean/} of
the accompanying repository, carry no \texttt{sorry}, and rest only on the standard axioms
\texttt{propext}, \texttt{Classical.choice} and \texttt{Quot.sound}. A claim described as
\emph{measured}, \emph{verified to} a stated tolerance, or \emph{checked numerically} is
computational evidence and is not a proof; the script producing it is named, and every named
script is run and its output compared against a recorded baseline
(\texttt{code/cite\_check.py}). A claim called a \emph{conjecture} is unproved, and is called
that even where the evidence is strong.



\section{Introduction and statement}

The expected characteristic polynomial of a random lift is the object at the
heart of the interlacing-families method of Marcus--Spielman--Srivastava
\cite{MSS}, who used the $r=2$ (signing) case to prove that bipartite Ramanujan
graphs exist in every degree; Hall--Puder--Sawin \cite{HPS} extended the
framework to arbitrary coverings. We evaluate it in closed form for the cycle.

A \emph{permutation $r$-lift} of a graph $G$ replaces each vertex by $r$
copies and each edge $\{u,v\}$ by the perfect matching $u_i\sim v_{\sigma(i)}$
of a permutation $\sigma\in S_r$ chosen independently and uniformly per edge.
For the cycle $C_n$, gauge-fixing a spanning path reduces the lift to a single
permutation $\sigma\in S_r$ on the closing edge, and the lift decomposes into
disjoint cycles: each $\ell$-cycle of $\sigma$ contributes a cycle $C_{n\ell}$
(fixed points contribute copies of $C_n$). Since
$\chi_{C_m}(x)=\det(xI-A_{C_m})=2T_m(x/2)-2$, the expected characteristic
polynomial is the expected cycle weight
\begin{equation}\label{eq:avg}
\Phi_{n,r}(x)\;=\;\E_{\sigma\in S_r}\;
\prod_{\ell\,\in\,\mathrm{cycles}(\sigma)}\bigl(2T_{n\ell}(x/2)-2\bigr).
\end{equation}

\begin{theorem}\label{thm:main}
For all $n\ge 1$, $r\ge 1$, with $Y=T_n(x/2)$,
\[
\Phi_{n,r}(x)\;=\;U_r(Y)-2U_{r-1}(Y)+U_{r-2}(Y)
\;=\;\underbrace{\bigl(2T_n(x/2)-2\bigr)}_{\chi_{C_n}(x)}\cdot\,
\underbrace{U_{r-1}\bigl(T_n(x/2)\bigr)}_{\Psi_r(x)} .
\]
\end{theorem}

The factored form identifies the quotient $\Psi_r$, by the theorem of
Hall--Puder--Sawin \cite{HPS} equating the expected new-eigenvalue
characteristic polynomial with the $d$-matching polynomial ($d=r-1$), as
\[
\mu_{d,C_n}(x)\;=\;U_{d}\bigl(T_n(x/2)\bigr),
\]
which is equivalent (through $\mathcal P_d(2t)=U_d(t)$, $\mathcal
C_n(2t)=2T_n(t)$ and the composition identity
$U_{dn+n-1}=U_d(T_n)\,U_{n-1}$) to \emph{Hall's conjecture}
$\mathcal C_{n,d}=\mathcal P_d(\mathcal C_n)$, proved combinatorially in \cite{CGHRS}.
That preprint has remained unrefereed since 2018 and carries no journal reference or
DOI as of August 2026. The reason is not mathematical. Chris Hall, who posed the
conjecture, has told us (personal communication, August 2026) that Andrew Herring was his
first doctoral student; Hall gave him the conjecture, Herring presented it as a problem at
a workshop, and the group of \cite{CGHRS} solved it there. They submitted the result, met
a couple of rejections, and abandoned the attempt for want of motivation rather than for
any defect. Hall regards the proof as correct, and so do we. We record the episode because
a reader weighing what is established about this identity should know both that the
combinatorial proof stands and that the present derivation is, as far as we can tell, the
first refereed or machine-checked one.
Our route is different and elementary; we make no use of matching-polynomial
combinatorics.

\begin{corollary}\label{cor:gf}
$\displaystyle\sum_{d\ge 0}\mu_{d,C_n}(x)\,z^d=\frac{1}{1-2T_n(x/2)\,z+z^2}$,
\quad and \quad
$\displaystyle\sum_{r\ge 0}\Phi_{n,r}(x)\,z^r=\frac{(1-z)^2}{1-2T_n(x/2)\,z+z^2}$.
\end{corollary}

\begin{corollary}[real-rootedness]\label{cor:roots}
All roots of $\Phi_{n,r}$ lie in $[-2,2]$: with $x=2\cos\theta$,
$\chi_{C_n}$ vanishes at $\theta=2\pi m/n$ and
$\Psi_r=\sin(rn\theta)/\sin(n\theta)$ vanishes at $\theta=m\pi/(rn)$ with $r\nmid m$
(when $r\mid m$ the denominator vanishes too and $\Psi_r=\pm r\neq0$); this leaves
exactly $n(r-1)$ roots, matching $\deg\Psi_r$.
\end{corollary}

\begin{corollary}[Fibonacci--Lucas evaluation]\label{cor:fib}
$\Phi_{n,r}(3)=(L_{2n}-2)\,F_{2nr}/F_{2n}$, since $T_n(3/2)=L_{2n}/2$ and
$U_{r-1}(L_{2n}/2)=F_{2nr}/F_{2n}$.
\end{corollary}

\begin{corollary}[parity]\label{cor:parity}
$\Psi_r(-x)=(-1)^{n(r-1)}\Psi_r(x)$: the quotient is supported on exponents
of a single parity, like a matching polynomial. (For bipartite $C_n$ it is an
even polynomial for every $r$; at $r=2$ it is the classical matching
polynomial $2T_n(x/2)$ of $C_n$, consistent with Godsil--Gutman
\cite{GG}.)
\end{corollary}

\subsection*{What this paper does}

The closed form is the starting point rather than the subject. Read as a statement about where the
roots of $\mu_{d,G}$ sit, it says that for the cycle they land exactly in
$\operatorname{spec}(T)$, and not merely in the interval $[-\rho,\rho]$ that Hall--Puder--Sawin
guarantee. The two differ whenever $T$ has a spectral gap, and whether the stronger statement holds
in general is the question here.

\S\ref{sec:uni} proves it for every graph of first Betti number one and every $d$.
\S\ref{sec:conj} states it as a conjecture for all graphs. \S\ref{sec:hallce} is the counterexample,
due to Chris Hall, together with the reason no local repair of it can work: by the Divisibility
Lemma of \S\ref{sec:divis} the branch polynomial divides $\mu_G$ across a separator, so the
offending root belongs to the divisor, and hypotheses on $G$ cannot see it. Minimum degree three
fails on $14$ vertices and $3$-connectivity on $23$. Supplementary material, including two families
in which the conjecture is provable and the record of the searches that missed the counterexample,
is in \cite{Supp}; the biregular case, which survives the refutation, is \cite{PartII}.

\section{Proof}

Write $f(\ell)=2T_\ell(Y)-2$ with $Y=T_n(x/2)$; by the composition identity
$T_{n\ell}=T_\ell\circ T_n$ this is the per-cycle factor in \eqref{eq:avg}.

\emph{Step 1 (exponential formula).} For any commutative ring $R$ and
$f:\mathbb N\to R$, the $S_r$-average
$\Phi_r=\E_\sigma\prod_{\ell\in\mathrm{cycles}(\sigma)}f(\ell)$ satisfies
\begin{equation}\label{eq:newton}
r\,\Phi_r=\sum_{k=1}^{r}f(k)\,\Phi_{r-k},\qquad\Phi_0=1 ,
\end{equation}
the coefficient form of $\sum_r\Phi_r z^r=\exp\bigl(\sum_\ell f(\ell)z^\ell/\ell\bigr)$.
(Proof: group $\sigma$ by cycle type $\lambda$ with Cauchy's count
$r!/z_\lambda$, mark a part, and re-index; see the Lean file
\texttt{Paper2ExpFormula.lean} for the fully detailed version.)

\emph{Step 2 (Chebyshev generating function).} With $f(\ell)=2T_\ell(Y)-2$,
\[
\sum_{\ell\ge1}f(\ell)\frac{z^\ell}{\ell}
=-\log\bigl(1-2Yz+z^2\bigr)+2\log(1-z)
\quad\Longrightarrow\quad
\sum_{r\ge0}\Phi_r z^r=\frac{(1-z)^2}{1-2Yz+z^2}.
\]
Extracting coefficients against $\sum U_r(Y)z^r=(1-2Yz+z^2)^{-1}$ gives
$\Phi_r=U_r-2U_{r-1}+U_{r-2}$, and the three-term recurrence
$U_r+U_{r-2}=2YU_{r-1}$ collapses this to $(2Y-2)\,U_{r-1}(Y)$.
\hfill$\square$

In the formalization the analytic Step 2 is replaced by an equivalent
finite argument: the closed form is shown to satisfy \eqref{eq:newton}
directly (a second-difference induction), and \eqref{eq:newton} with
$\Phi_0=1$ determines $\Phi$ over a characteristic-zero domain.

\section{Graphs of first Betti number one}\label{sec:uni}

The proof above uses cycles in exactly one place: gauge-fixing a spanning path
leaves a \emph{single} permutation, because $C_n$ has $|E|-|V|+1=1$. The
argument therefore applies verbatim to any connected $G$ with
$b_1(G)=|E|-|V|+1=1$, a cycle with trees attached, and we record what it gives.

By \cite{HPS} the expected characteristic polynomial of a random $r$-lift of
\emph{any} graph factors as $\chi_G\cdot\mu_{r-1,G}$; the content below is
therefore not that factorization but a closed form for the factor
$\mu_{d,G}$ itself. Question 6.5 of \cite{HPS} asks which parts of the theory
of the matching polynomial $\mu_{1,G}$ generalize to $\mu_{d,G}$; the following
answers it for $b_1(G)=1$.

Let $C$ be the unique cycle of $G$ and let $G-V(C)$ be the forest obtained by
deleting its vertices. Write $\mu_H$ for the ordinary matching polynomial of a
graph $H$, with $\mu_\emptyset=1$. Define
\begin{equation}\label{eq:V}
V_0=1,\qquad V_1=\mu_G,\qquad
V_d=\mu_G\,V_{d-1}-\mu_{G-V(C)}^{\,2}\,V_{d-2}.
\end{equation}

\begin{theorem}\label{thm:uni}
For every connected $G$ with $b_1(G)=1$ and every $d\ge0$, $\mu_{d,G}=V_d$.
\end{theorem}

Theorem~\ref{thm:uni} is proved below, through the twisted characteristic polynomial; it is not
formalized.

For $G=C_n$ the deleted forest is empty, so $\mu_{G-V(C)}=1$ and \eqref{eq:V}
is the recurrence for $U_d(\mu_{C_n}/2)=U_d(T_n(x/2))$: Theorem~\ref{thm:main}
is the case $\mu_{G-V(C)}=1$. In general the attached trees enter only through
$\mu_{G-V(C)}$; for the triangle with one pendant edge,
$\mu_G=x^4-4x^2+1$ and $\mu_{G-V(C)}=x$.

The proof goes through the twisted characteristic polynomial. Let $e=uv$ be the
unique non-tree edge, let $A_G(z)$ be the adjacency matrix of $G$ with the
$(u,v)$ entry replaced by $z$ and the $(v,u)$ entry by $z^{-1}$, and set
$F(x,z)=\det(xI-A_G(z))$.

\begin{lemma}\label{lem:PQ}
$F(x,z)=\mu_G(x)-\mu_{G-V(C)}(x)\,(z+z^{-1})$.
\end{lemma}

\begin{proof}
By the Sachs coefficient theorem $\det(xI-A)$ expands over elementary
subgraphs, that is disjoint unions of edges and cycles. Since a permutation
uses each of the two twisted entries at most once and
$A_G(z)^{\mathsf T}=A_G(z^{-1})$, the determinant is of the form
$P(x)+Q(x)(z+z^{-1})$. Replacing $z$ by $\pm1$ and averaging kills exactly the
elementary subgraphs whose cycle traverses $e$; as $G$ has the single cycle $C$,
those are the subgraphs containing $C$, so $P=\mu_G$, and the surviving terms
contribute $Q=-\mu_{G-V(C)}$, a cycle component carrying the Sachs weight $-2$.
\end{proof}

\begin{proof}[Proof of Theorem~\ref{thm:uni}]
Gauge-fix a spanning tree: the $r$-lift is determined by one permutation
$\sigma\in S_r$ on $e$, and its components are indexed by the cycles of
$\sigma$, the component of an $\ell$-cycle being the $\ell$-fold cyclic cover
$G_\ell$. Decomposing the $\mathbb{Z}/\ell$ action into characters gives the
classical factorization $\chi_{G_\ell}(x)=\prod_{j=0}^{\ell-1}F(x,\omega^j)$
with $\omega=e^{2\pi i/\ell}$ \cite{MS,DFMRS}. With $P,Q$ as in
Lemma~\ref{lem:PQ} and $y=-P/(2Q)$,
\[
\chi_{G_\ell}=\prod_{j=0}^{\ell-1}\Bigl(P+2Q\cos\tfrac{2\pi j}{\ell}\Bigr)
=(-Q)^{\ell}\bigl(2T_\ell(y)-2\bigr),
\]
so by the exponential formula \eqref{eq:newton}, with $w=-Qz$,
\[
\sum_{r\ge0}\Phi_{G,r}z^{r}
=\exp\Bigl(\sum_{\ell\ge1}\bigl(2T_\ell(y)-2\bigr)\frac{w^{\ell}}{\ell}\Bigr)
=\frac{(1-w)^2}{1-2yw+w^2},
\]
and extracting coefficients as in Section~2 gives
$\Phi_{G,r}=\chi_G\cdot(-Q)^{r-1}U_{r-1}(y)$. Comparing with
$\Phi_{G,r}=\chi_G\cdot\mu_{r-1,G}$ \cite{HPS} yields
$\mu_{d,G}=(-Q)^{d}U_d(y)$, which satisfies \eqref{eq:V}.
\end{proof}

Since $U_d(y)=2^{d}\prod_{k=1}^{d}(y-\cos\tfrac{k\pi}{d+1})$, the theorem can
be restated as a product,
$\mu_{d,G}(x)=\prod_{k=1}^{d}F\bigl(x,e^{\,ik\pi/(d+1)}\bigr)$, whence:

\begin{corollary}\label{cor:unirr}
For $b_1(G)=1$, $\mu_{d,G}$ is a product of $d$ real-rooted polynomials, and its
roots are exactly the eigenvalues of the Hermitian matrices
$A_G\bigl(e^{\,ik\pi/(d+1)}\bigr)$, $1\le k\le d$.
\end{corollary}

This is the real-rootedness theorem of \cite{HPS} for this family, with the
roots identified. The localization itself, that a vanishing $V_d$ forces
$|\mu_G|\le2|\mu_{G-V(C)}|$, so that a root can never sit in a spectral gap, is machine-checked in Lean~4 as \texttt{cV\_root\_mem\_band}, resting on
\texttt{cU\_ne\_zero\_of\_one\_le\_abs}: the Chebyshev polynomial $U_d$ has no
root of absolute value at least one.

The roots can be located much more precisely than that. Since $b_1(G)=1$ the
group $\pi_1(G)$ is $\mathbb{Z}$, so the universal cover $T$ of $G$ \emph{is}
the $\mathbb{Z}$-cover, and $T$ is a periodic decorated chain whose adjacency
operator decomposes by Floquet--Bloch into the fibres $A_G(e^{i\theta})$. Hence
\[
\operatorname{spec}(T)=\bigcup_{\theta}\operatorname{spec}A_G(e^{i\theta})
=\Bigl\{x:\ \exists\,s\in[-1,1],\ \mu_G(x)=2s\,\mu_{G-V(C)}(x)\Bigr\},
\]
a finite union of bands, together with a flat band at each common zero of
$\mu_G$ and $\mu_{G-V(C)}$, where $F(x,\cdot)$ vanishes identically.

\begin{corollary}\label{cor:floquet}
For $b_1(G)=1$ the roots of $\mu_{d,G}$ are exactly the solutions of
$\mu_G(x)=2\cos\!\bigl(\tfrac{k\pi}{d+1}\bigr)\mu_{G-V(C)}(x)$ for
$1\le k\le d$. In particular every root of every $\mu_{d,G}$ lies in
$\operatorname{spec}(T)$, and as $d\to\infty$ the roots equidistribute on
$\operatorname{spec}(T)$ with respect to its density of states.
\end{corollary}

\begin{proof}
Immediate from Theorem~\ref{thm:uni} in the form
$\mu_{d,G}=\prod_{k}F(x,e^{ik\pi/(d+1)})$ and
$F(x,e^{i\theta})=\mu_G-2\cos\theta\,\mu_{G-V(C)}$. The equidistribution is the
arcsine law for $\cos\bigl(k\pi/(d+1)\bigr)$ transported through the fibre map.
\end{proof}

Hall--Puder--Sawin place the roots of $\mu_{d,G}$ in the Ramanujan
\emph{interval} $[-\rho,\rho]$, $\rho$ being the spectral radius of $T$.
Corollary~\ref{cor:floquet} places them in $\operatorname{spec}(T)$, which is a
union of bands and is a proper subset of that interval whenever $T$ has a
spectral gap: no root of any $\mu_{d,G}$ can ever enter a gap. The two
statements coincide exactly when the band structure is trivial, and that is
precisely the case of a cycle, where $\mu_{G-V(C)}=1$ and
$\operatorname{spec}(T)=[-2,2]$ is the whole interval. Attaching even a single
pendant edge opens gaps, and they are not small (computed band structures,
proportion of $[-\rho,\rho]$ that is a gap):
\[
\begin{array}{lcc}
\hline
G & \rho & \text{gap proportion}\\
\hline
C_n & 2 & 0\%\\
\text{triangle}+\text{pendant} & 2.170 & 36.5\%\\
C_4+\text{pendant} & 2.136 & 45.9\%\\
\text{triangle}+P_2 & 2.214 & 51.3\%\\
C_4+P_3 & 2.189 & 54.1\%\\
\text{triangle}+K_{1,3} & 2.514 & 65.9\%\\
\hline
\end{array}
\]
So the cycle is not merely the easiest case of Theorem~\ref{thm:uni}: it is the
degenerate one, the unique member of the family whose spectrum has no gaps and
for which the Ramanujan interval is attained. Any cycle with a tree attached
exhibits a band structure that no cycle can display.

\paragraph{The boundary is sharp.}
For $b_1(G)=b$ the lift is determined by $b$ permutations, that is by an action
of the free group $F_b$ on $[r]$, whose orbits are the components of the lift.
The exponential formula still applies, in the form
$N_r=\sum_{\ell}\binom{r-1}{\ell-1}C_\ell N_{r-\ell}$ with
$N_r=(r!)^{b}\Phi_{G,r}$ and $C_\ell$ the sum of $\chi$ over the transitive
$b$-tuples on $[\ell]$; we have verified this for $b\le3$. What fails is the
atom. For $b=1$ every transitive tuple gives the \emph{same} cover $G_\ell$, so
$C_\ell=(\ell-1)!\,\chi_{G_\ell}$: one graph per $\ell$, with $\chi_{G_\ell}$
Chebyshev in $\ell$. For $b\ge2$ the transitive tuples number $(\ell-1)!$ times
the count of index-$\ell$ subgroups of $F_b$ and spread over many
non-isomorphic covers with distinct spectra, so $C_\ell$ carries no such
structure. Accordingly Theorem~\ref{thm:uni} is not merely unproved for
$b\ge2$: it is false there already at $d=1$, as we checked for the theta graph,
the bowtie, $K_4$, $K_{2,3}$ and two triangles joined by an edge.

Theorem~\ref{thm:uni} was verified by exact integer computation against the
definition of $\mu_{d,G}$ as an average of matching polynomials over all
$d$-covers, and against brute-forced lift averages, on thirteen named unicyclic
graphs for $r\le6$ and on a random sample of thirty-four for $r\le5$;
Lemma~\ref{lem:PQ} was checked on eleven.

\section{A conjecture for all graphs}\label{sec:conj}
Corollary~\ref{cor:floquet} is proved by Floquet theory, which is available
only because $\pi_1(G)=\mathbb{Z}$. Its \emph{statement}, however, makes sense
for every graph, and we know one case of it: at $b_1(G)=1$ it is
Corollary~\ref{cor:floquet}, for every $d$. We therefore propose:

\begin{conjecture}\label{conj:gap}
For every finite graph $G$ and every $d\ge1$, all roots of $\mu_{d,G}$ lie in
$\operatorname{spec}(T)$, the spectrum of the universal cover of $G$.
\end{conjecture}

\begin{theorem}[Hall \cite{HallCE}]\label{thm:hallce}
Conjecture~\ref{conj:gap} is false, already at $d=1$ and already for simple connected
loopless bipartite graphs.
\end{theorem}

The counterexample and its certificate are entirely due to Chris Hall. We record it here
because it settles the central question of this paper, and because everything else in the
note has to be read in its light. \S\ref{sec:hallce} states his construction, our
independent verification of it, and what we can add.

The biregular case of Conjecture~\ref{conj:gap} at $d=1$ is not new. It is
Problem~1 of Song, Fan and Miao \cite{SFM}, posed in 2023: they ask whether the
matching polynomial of the incidence graph $B_H$ of a $(d,r)$-regular hypergraph
has no root $\mu$ with $0<|\mu|<|\sqrt{d-1}-\sqrt{r-1}|$, and every
$(d,r)$-biregular bipartite graph is such an incidence graph. Their Theorem 1.8
attaches a consequence: an affirmative answer would give every $(d,r)$-regular
hypergraph a left-sided Ramanujan $2$-covering. The only bound in that direction
we know of is their Lemma 5.5, which goes the other way,
$0<\mu_{\tau}(B_H)\le\max\{\sqrt d,\sqrt r\}$. We are not aware of any proof,
partial result or counterexample. Conjecture~\ref{conj:gap} extends their problem
in two directions, to arbitrary graphs and to all $d$.

Question 6.3 of \cite{HPS} asks the corresponding question for graph
eigenvalues, and leaves it open: it distinguishes graphs whose nontrivial
spectrum lies in the \emph{interval} $[-\rho,\rho]$ from those whose nontrivial
spectrum lies in $\operatorname{spec}(T)$, observes that the two notions
separate exactly when $T$ has a gap, and gives the biregular tree as the
example. Conjecture~\ref{conj:gap} is the matching-polynomial form of that
distinction.

The conjecture strengthens the root-localization theorem of \cite{HPS}, which
places the roots in the interval $[-\rho,\rho]$ with $\rho$ the spectral radius
of $T$;
the conjecture places them in $\operatorname{spec}(T)$ itself. The two agree
whenever $T$ has no spectral gap, in particular for regular $G$, where $T$
is a regular tree and $\operatorname{spec}(T)=[-2\sqrt{q},2\sqrt{q}]$ is an
interval. They differ exactly when $T$ has a gap, and the smallest examples are
biregular: for $K_{2,3}$ the universal cover is the $(3,2)$-biregular tree,
with spectrum $\{0\}\cup\pm[\sqrt2-1,\sqrt2+1]$ and hence a gap
$(0,\sqrt2-1)$ inside the Ramanujan interval. We have tested the conjecture on the
biregular family, where the universal cover is an $(a,b)$-biregular tree with
spectrum $\{0\}\cup\pm[\,|s-t|,s+t\,]$, $s=\sqrt{a-1}$, $t=\sqrt{b-1}$, so the
gap is $(0,|s-t|)$. In each case $\mu_{d,G}$ was computed by averaging matching
polynomials over all $d$-covers. No root ever lies in the gap; we also record
the smallest nonzero $|{\rm root}|$ as a multiple of the gap edge, which
measures how close the bound comes to being attained.
\[
\begin{array}{lccccc}
\hline
G & b_1 & \text{type} & \text{gap edge} & d & \min|{\rm root}|/\text{gap edge}\\
\hline
K_{2,3} & 2 & (3,2) & 0.4142 & 1,2,3,4 & 2.72,\ 2.19,\ 1.94,\ 1.78\\
K_{2,4} & 3 & (4,2) & 0.7321 & 1,2,3 & 1.93,\ 1.65,\ 1.52\\
K_{2,5} & 4 & (5,2) & 1.0000 & 1,2,3 & 1.66,\ 1.46,\ 1.37\\
K_{3,4} & 6 & (4,3) & 0.3178 & 1,2 & 3.04,\ 2.31\\
K_{3,5} & 8 & (5,3) & 0.5858 & 1,2 & 2.10,\ 1.73\\
\text{subdivision of }K_4 & 3 & (3,2) & 0.4142 & 1,2 & 1.97,\ 1.67\\
\hline
\end{array}
\]
All of these have $b_1\ge2$, so none is covered by
Corollary~\ref{cor:floquet}. It is worth recording that even $d=1$ is not
formal. There $\mu_{1,G}$ is the ordinary matching polynomial, whose roots are
eigenvalues of the path tree of $G$ (Godsil), and the path tree embeds in $T$;
but eigenvalues of a finite subtree of $T$ need not lie in
$\operatorname{spec}(T)$ when $T$ has a gap. For the $(3,2)$-biregular tree,
whose gap is $(0,\sqrt2-1)$, the path $P_8$ is a subtree with eigenvalue
$2\cos(4\pi/9)=0.3473$ inside that gap. So the $d=1$ case does not follow from
the path-tree argument. It is nevertheless a theorem for an infinite family.

\section{Hall's counterexample}\label{sec:hallce}

Everything in this section is due to Chris Hall, who
constructed the graph, found the certificate, and wrote the verification note
\cite{HallCE}. Our contribution is the independent check, the formalization of the
algebra. The two-parameter family containing his graph is in \cite{Supp}.

\subsection{The graph and its matching polynomial}

Let $G$ consist of a central vertex $c$; for $1\le i\le5$ a copy of $K_{2,5}$ with
degree-five vertices $v_i,w_i$ and middle vertices $u_{i,1},\dots,u_{i,5}$; a pendant leaf
$\ell_i$ at each $w_i$; and the edges $cv_i$. It is simple, connected, loopless and
bipartite, with $41$ vertices, $60$ edges and $b_1(G)=20$.

Writing $H$ for one branch after deleting $c$, the vertex-deletion recurrence gives
$\mu_{H-v}=x^5(x^2-6)$ and $\mu_H=x^4(x^4-11x^2+25)$, and then
\begin{equation}\label{eq:hallmu}
   \mu_G(x)=\mu_H(x)^4\bigl(x\mu_H(x)-5\mu_{H-v}(x)\bigr)
   =x^{21}(x^4-11x^2+25)^4(x^2-5)(x^2-11),
\end{equation}
so $\sqrt5$ is a simple root. Identity \eqref{eq:hallmu}, the simplicity of the root, and
the value of the cofactor are machine-checked in \texttt{RamaLean/HallCounterexample.lean}
as an identity in $\mathbb Z[X]$ (\texttt{muG\_eq}, \texttt{muG\_root\_sqrt5},
\texttt{muG\_root\_simple}).

\subsection{The spectral exclusion}

Angel, Friedman and Hoory \cite{AFH} characterise invertibility of $A_T-\lambda I$ by the
existence of a finite nonzero \emph{ratio system}: an assignment $r_e$ to directed edges
with $\lambda=1/r_e+\sum_{e\to f}r_f$, whose decay matrix $\mathcal R_{e,f}=|r_f|^2$ on
non-backtracking pairs has spectral radius below one. Hall exhibits such a system at
$\lambda=\sqrt5$ with eight edge types and ratios in $\mathbb Q(\sqrt5,\sqrt{41})$. Its
decay matrix has $120$ states; the ten states at the leaves are transient, and the
recurrent block of $110$ states reduces by symmetry to a $6\times6$ orbit quotient $K$ for
which
\[
   x=(20,121,47,33,28,16)^{\mathsf T}\quad\text{satisfies}\quad Kx<x
\]
componentwise, so $\rho(K)<1$ by Perron--Frobenius. Hence $\sqrt5\notin
\operatorname{spec}(A_T)$, and with \eqref{eq:hallmu} this proves
Theorem~\ref{thm:hallce}. Since $G$ is bipartite the same holds at $-\sqrt5$, and
Hall--Puder--Sawin place $\pm\sqrt5$ inside $[-\rho(T),\rho(T)]$, so these are roots in a
genuine internal gap.

The eight ratio identities, the elementary bounds $32/5<\sqrt{41}<13/2$, the strict
positivity of all six components of $x-Kx$, and a Collatz--Wielandt lemma (a nonnegative
matrix with a positive strictly subinvariant vector has spectral radius below one) are
machine-checked in \texttt{RamaLean/HallCounterexample.lean}. The argument common to all three
counterexamples of this paper is isolated once in
\texttt{RamaLean/RatioCertificate.lean}: a ratio system with its decay matrix
(\texttt{decay}, \texttt{decay\_nonneg}), the Collatz--Wielandt step in the strict form the
certificate needs (\texttt{decay\_lt\_one}), and the assembly with Angel--Friedman--Hoory as an
explicit hypothesis (\texttt{spec\_excluded}). The three instances differ only in the numbers
supplied. A schema rather than closed-form algebra is the right level here: Hall's ratios lie
in $\mathbb{Q}(\sqrt5,\sqrt{41})$, but that is a feature of his graph, and for the $31$-vertex
one the ratios admit no minimal polynomial of small degree and moderate height. The theorem of \cite{AFH} is
not in Mathlib and enters as an explicit hypothesis, in the shape it is consumed.

We verified the certificate independently, reconstructing every object from the graph
rather than copying the displayed matrices: the matching numbers
$(1,60,1585,24260,238105,1564976,6973855,20805500,39784375,44062500,21484375)$; the ratio
equations on all $120$ directed edges; the strongly connected components of the decay
digraph; the orbit quotient rebuilt from follower counts; and the decay rate by two
independent routes, the componentwise certificate and the characteristic polynomial
$t^6-\tfrac{201+19s}{800}t^2+\tfrac{21s-241}{250}$ with $s=\sqrt{41}$. That reconstruction is
\texttt{code/hall\_certificate.py}, which rebuilds each object from the construction and copies no
displayed matrix: the matching numbers agree term for term and $\mu_G(\sqrt5)$ vanishes to
$10^{-45}$; the ratio system reduces to eight unknowns by symmetry and is solved exactly in
$\mathbb Q(\sqrt5,\sqrt{41})$, its solution satisfying all $120$ equations to $4\times10^{-16}$;
the ten transient states are exactly the leaf-incident directed edges and the recurrent block has
$110$; and the six components of $x-Kx$ reproduce the exact expressions carried by
\texttt{RamaLean/HallCounterexample.lean} to $3\times10^{-14}$, with $\rho(K)=0.9636233789$. Three
of the eight ratios are negative, which is why a positive cavity iteration does not find the
system. As a physical check,
the non-backtracking cavity Green function at $\sqrt5+i\eta$ converges to Hall's real
ratios as $\eta\downarrow0$, with density of states vanishing linearly in $\eta$.

\subsection{The mechanism is a separation}\label{sec:mechanism}

Every counterexample is built on a separation. If a vertex $v$ separates $G$ into $p$ isomorphic
branches $H$, the deletion recurrence gives
\begin{equation}\label{eq:branch}
   \mu_G=\mu_H^{\,p-1}\bigl(x\mu_H-p\,\mu_{H-v}\bigr),
\end{equation}
and the root is placed by tuning $p$ and $H$; this is
\texttt{CutVertexMechanism.branch\_factor}, with \texttt{root\_of\_branch\_factor} transferring
a root of the last factor to $\mu_G$. Hall's graph and the graphs of
\cite{Supp} are of this shape.

A cut vertex is only the smallest separation. Two hubs $u,v$ with $p$ identical branches
attached to both leave no cut vertex, since removing either hub leaves the graph joined through
the other, while $\{u,v\}$ separates it into $p$ pieces. Writing $A=\mu_B$ for the branch,
$A_1,A_2$ for the branch with one anchor deleted and $A_{12}$ for both, deletion at $u$ and then
at $v$ gives
\begin{equation}\label{eq:twocut}
   \mu_G=A^{\,p-2}\bigl(x^2A^2-px(A_1+A_2)A+pA_{12}A+p(p-1)A_1A_2\bigr),
\end{equation}
which is \texttt{SeparationOrder.two\_cut\_factor}, with \texttt{root\_of\_two\_cut}.

The bracket in \eqref{eq:twocut} is \emph{quadratic} in $p$ where the bracket in
\eqref{eq:branch} is linear. Read as polynomials in the branch count the two have degrees one and
two, with leading coefficient $A_1A_2$ in the second case
(\texttt{cutBracket\_natDegree}, \texttt{twoCutBracket\_natDegree}, \texttt{freedom\_grows}).
Excluding cut vertices does not remove the engine; it upgrades it.

\begin{proposition}\label{prop:twoconn}
There are $2$-connected counterexamples to Conjecture~\ref{conj:gap}.
\end{proposition}

Let $B_{q,T}$ be $K_{2,q}$ with a path of $T$ further vertices appended to one of its two
degree-$q$ vertices; attach the other degree-$q$ vertex to $u$ and the end of the path to $v$,
and take $p$ copies.

\begin{center}
\begin{tabular}{lrrrrrr}
branch & $p$ & $n$ & $\kappa(G)$ & $\theta$ & defect & gap width\\\hline
$B_{4,2}$ & $7$ & $58$ & $2$ & $2.137920$ & $0.003574$ & $0.14758$\\
$B_{5,2}$ & $6$ & $56$ & $2$ & $2.287307$ & $0.005624$ & $0.06636$\\
$B_{5,3}$ & $6$ & $62$ & $2$ & $2.317609$ & $0.008332$ & $0.09876$\\
\end{tabular}
\end{center}

All three are simple, bipartite, of minimum degree two and maximum degree at most seven, and of
vertex connectivity exactly two, computed by Menger as a unit-capacity maximum flow over every
non-adjacent pair. Each $\theta$ is an algebraic number of degree $12$ or $14$ isolated from the
exact integer bracket, and \eqref{eq:twocut} is checked against a brute-force matching
polynomial. That $\theta\notin\operatorname{spec}(T)$ is confirmed twice: the
Angel--Friedman--Hoory decay rate at $40$ digits is $0.943$ to $0.959$, and the density of states
of the universal cover, from the complex cavity equations, falls by a factor of ten per decade of
$\eta$ at $\theta$ while an in-band control holds constant to four figures
(\texttt{code/twocut.py}).

Equation \eqref{eq:twocut} is the case $k=2$ of a general expansion. With $k$ hubs and $p$
branches attached to all of them,
\begin{equation}\label{eq:kcut}
   \mu_G=\sum_{S\subseteq[k]}(-1)^{|S|}x^{\,k-|S|}
        \sum_{P\vdash S}(p)_{|P|}\Bigl(\prod_{J\in P}A_J\Bigr)A^{\,p-|P|},
\end{equation}
with $A_J$ equal to $\mu_B$ with the anchors indexed by $J$ deleted and $(p)_m$ the falling
factorial: $A^{p-k}$ times a bracket of degree $k$ in $p$. A $k$-cut thus needs only
$k$-connectedness while supplying a degree-$k$ polynomial to tune. No counterexample above
connectivity two is known: a search over thirty-eight branch families at $k=3,4,5$ with up to
$62$ vertices, with \eqref{eq:kcut} checked against brute force at $k=3$, produced none, and it
is not vacuous, those graphs carrying two to five gaps of total width $0.40$ to $0.75$
(\texttt{code/kcut.py}).

\subsection{The Divisibility Lemma}\label{sec:divis}

The counterexamples are all instances of one algebraic fact, which we state once.

\begin{lemma}[Divisibility]\label{lem:divis}
Let $G$ be a finite graph, $S\subseteq V(G)$ with $|S|=k$, and suppose $G-S$ has $p$ components,
each inducing a copy of the same graph $B$. Write $A=\mu_B$. Then
\[
  A^{\,p-k}\ \big|\ \mu_G .
\]
\end{lemma}

\begin{proof}
Expand $\mu_G=\sum_M(-1)^{|M|}x^{n-2|M|}$ over matchings $M$ of $G$ and group the terms by
\[
  J(M)=\{\,i : \text{some edge of $M$ joins $S$ to the $i$-th component}\,\}.
\]
A matching uses each vertex of $S$ at most once, so distinct components in $J(M)$ are entered
through distinct vertices of $S$; choosing one such vertex per component is an injection
$J(M)\hookrightarrow S$, whence $|J(M)|\le k$. For $i\notin J(M)$ the restriction of $M$ to the
$i$-th component is an arbitrary matching of $B$, and summing over those choices independently
contributes a factor $A^{\,p-|J(M)|}$. Each term of the resulting sum is therefore divisible by
$A^{\,p-k}$, and hence so is $\mu_G$.
\end{proof}

At $k=2$ this is the divisor $A^{p-2}$ of the two-cut identity. The exponent is positive exactly
when $p>k$, which is why every construction below needs at least $k+1$ branches. The counting step
and the algebraic step are machine-checked separately
(\texttt{RamaLean/DivisibilityLemma.lean}: \texttt{card\_branches\_le\_sep},
\texttt{dvd\_sum\_of\_pow\_terms}, \texttt{matching\_expansion\_dvd}); the matching expansion
itself is not, Mathlib carrying no matching polynomial.

\subsection{Neither minimum degree nor connectivity}\label{sec:mindeg3}

Connectivity is not the right invariant, because a graph of vertex connectivity $\kappa$ has
minimum degree at least $\kappa$, and it is the minimum degree that the counterexamples respond
to. Every counterexample above has minimum degree at most two: Hall's has leaves, the
$31$-vertex graph of \cite{Supp} has degree-2 vertices, and so do all three
graphs of Proposition~\ref{prop:twoconn}.

The two-hub construction separates the two conditions exactly, since it can be run with branches
whose every vertex has degree at least three: ladders anchored at both ends of both rails,
prisms, and $K_{3,q}$ blocks with two of the three degree-$q$ vertices anchored. That yields
$102$ graphs of minimum degree three whose vertex connectivity is \emph{two}, so each retains the
separating pair the engine needs, and $77$ of them have a gap of width at least $0.05$, several
wider than $2.9$, against the $0.066$ to $0.148$ that the counterexamples of
Proposition~\ref{prop:twoconn} occupy. None has a root of the bracket in a gap
(\texttt{code/mindeg3.py}); that sweep did not test the roots of $A$, which is where the
counterexample below lives.

\begin{conjecture}[D3]\label{conj:mindeg3}
Every finite graph of minimum degree at least three satisfies
$\operatorname{Zeros}(\mu_G)\subseteq\operatorname{spec}(T_G)$.
\end{conjecture}

\textbf{Conjecture~\ref{conj:mindeg3} is false.} We record the counterexample, which is smaller
than Hall's by a factor of three.

\begin{proposition}[D3 fails on $14$ vertices]\label{prop:d3false}
Let $G$ have two hubs $u,v$, three copies of the star $K_{1,3}$ with centre $c_i$ and leaves
$\ell_{i,1},\ell_{i,2},\ell_{i,3}$, and every leaf joined to both hubs. Then $G$ is simple,
connected and bipartite on $14$ vertices and $27$ edges, with
$\deg c_i=\deg\ell_{i,j}=3$ and $\deg u=\deg v=9$, so $\delta(G)=3$; and
\[
  \mu_G=x^4\,(x^2-3)\,(x^8-24x^6+171x^4-396x^2+270),
\]
so $\sqrt3$ is a root of $\mu_G$, while $\sqrt3\notin\operatorname{spec}(T_G)$.
\end{proposition}

Both halves are exact. That $\mu_G$ factors as displayed is computed from the graph by the deletion
recursion, with $(x^2-3)$ divided out over $\mathbb{Q}$. That
$\sqrt3\notin\operatorname{spec}(T_G)$ is Proposition~\ref{prop:d3cert}, whose arithmetic is
machine-checked (\texttt{RamaLean/D3Counterexample.lean}); the analytic input is the
Angel--Friedman--Hoory criterion \cite{AFH}, which is cited and not reproved.

\begin{proposition}[the certificate]\label{prop:d3cert}
On the four directed-edge orbits of $G$, write $g_1$ for centre-to-leaf, $g_2$ for leaf-to-centre,
$g_3$ for leaf-to-hub and $g_4$ for hub-to-leaf. At $\lambda=\sqrt3$ the cavity system
\[
 \lambda=\tfrac1{g_1}+2g_3,\quad \lambda=\tfrac1{g_2}+2g_1,\quad
 \lambda=\tfrac1{g_3}+8g_4,\quad \lambda=\tfrac1{g_4}+g_2+g_3
\]
is solved by $g_1=1/\sqrt3$, $g_2=\sqrt3$, $g_3=0$, $g_4=\infty$, the last pair escaping along
$g_3g_4\to-\tfrac18$. The decay quotient has characteristic polynomial $t^4-\tfrac18t^2-\tfrac12$,
whose largest root in modulus is
\[
  \rho=\sqrt{\tfrac{1+\sqrt{129}}{16}}=\tfrac14\sqrt{1+\sqrt{129}}=0.878842\ldots<1,
\]
and the vertex Green's functions are $G_{\mathrm{hub}}=\sqrt3/3$,
$G_{\mathrm{centre}}=-\sqrt3/6$, $G_{\mathrm{leaf}}=0$, all finite.
\end{proposition}

Three things make this work where the naive attempt fails. The pole in $g_4$ is a pole of the
\emph{coordinate}: the decay quotient's entries diverge, but its characteristic polynomial does not,
because a determinant collects only vertex-disjoint cycle collections and the two cycles present,
$(g_3g_4)$ and $(g_1g_3g_4g_2)$, share vertices, so their weights $8(g_3g_4)^2=\tfrac18$ and
$32g_1^2g_2^2(g_3g_4)^2=\tfrac12$ are finite and there is no cross term. Next, $\rho<1$ needs no
surd: with $s=t^2$ the polynomial is $8s^2-s-4$, and $(s-1)^2\ge0$ gives $8s^2-16s+8\ge0$, which on
substituting $8s^2=s+4$ leaves $15s\le12$, so $s\le\tfrac45$
(\texttt{decay\_root\_lt\_one}). Last, the finiteness of every $G_w$ excludes the one remaining
possibility, that $\sqrt3$ is an \emph{isolated point} of $\operatorname{spec}(T_G)$: an atom forces
a pole of $G_w$ with the eigenvector weight as residue, and there is none. That step is not a
formality, since $\operatorname{spec}(T)$ is closed and the biregular cover of \S\ref{sec:conj}
carries exactly such an isolated $\{0\}$.

As an independent check, the exact $\sum_wG_w=\sqrt3/6$ is reproduced by the numerical resolvent to
$7.7\times10^{-10}$ at $\eta=10^{-6}$, the error falling as $\eta^2$, and the exact $\rho$ lies
between the decays $0.892472$ and $0.867159$ computed at $\sqrt3\mp0.005$
(\texttt{code/d3\_exact.py}).

The mechanism is Hall's verbatim. His violating root $\sqrt5$ is the top matching root of the
star $K_{1,5}$, whose matching polynomial is $x^4(x^2-5)$; the divisor is a star. A star's leaves
have degree one \emph{in the divisor} however large their degree is \emph{in $G$}, so a hypothesis
on $\delta(G)$ constrains the ambient graph and not the divisor. Here the branch is $K_{1,3}$,
$A=\mu_{K_{1,3}}=x^2(x^2-3)$, and the two-cut identity
$\mu_G=A^{p-2}\bigl(x^2A^2-px(B_u+B_v)A+pDA+p(p-1)B_uB_v\bigr)$ makes every root of $A$ a root of
$\mu_G$ once $p\ge3$. Three branches are needed and two do not suffice, since $A^0=1$ carries no
root.

This is exactly what the searches above could not see, and the reason is worth stating because it
is not a matter of scale. Each of them lifted the block's degrees to three by adding edges
\emph{inside} the block, which changes $A$ and moves its roots; the counterexample lifts them from
\emph{outside}, by attaching each leaf to both hubs, which leaves $A$ untouched. The five clean
searches were therefore all run on blocks whose matching polynomial had already been perturbed
away from the star. Separately, \texttt{code/mindeg3.py} iterated over the roots of the bracket
and never over the roots of $A$, so the divisibility that is the whole engine went untested; that
is repaired in \texttt{code/mindeg3adv.py}, which reports a signed margin rather than a verdict.

One trap on the way to Proposition~\ref{prop:d3cert} is worth recording, because it inverts the
answer rather than blurring it, and because the same trap is what let five searches come back
clean. Solved naively, the Angel--Friedman--Hoory system
$\lambda=1/r_e+\sum_{e\to f}r_f$ has four orbits here under $\operatorname{Aut}(G)$, and
eliminating gives a quadratic whose two real branches both have decay above one, which reads as
$\sqrt3\in\operatorname{spec}(T_G)$. That is wrong: a hair either side of $\sqrt3$ there are
\emph{three} real branches and the third has decay about $0.87$. Its $r(\ell\to u)$ crosses zero
exactly at $\sqrt3$, so the elimination multiplies by it and drops the branch, and a pointwise
solver returns nothing there for the same reason. Proposition~\ref{prop:d3cert} is what the branch
becomes once the singularity is removed rather than divided by.

Independently of the algebra, the density-of-states ladder settles the same point numerically and is
what located it in the first place: with $g_e=1/(z-\sum_{e\to f}g_f)$ at $z=\lambda+i\eta$, the
quantity $|\operatorname{Im}\sum_w G_w|$ decays linearly in $\eta$ outside the spectrum, tends to a
positive constant inside a band, and diverges like $\eta^{-1}$ at an atom. On this graph it reports
an atom at $0$, a band at $1.68$, a band at $2.40$, and linear decay at $\sqrt3$, the atom at the
origin serving as the internal control (\texttt{code/d3\_counterexample.py}).

Proposition~\ref{prop:d3false} does not touch the connectivity statement, since $G$ there has vertex
connectivity two. But C2 falls to the same construction, and the reason it does is the point of this
subsection.

\begin{proposition}[C2 fails on $23$ vertices]\label{prop:c2false}
Let $G$ have $k$ hubs, $p$ copies of $K_{1,m}$, and every leaf joined to \emph{all} $k$ hubs. A
matching sends each hub into at most one branch, so at least $p-k$ branches contribute their full
matching polynomial and
\[
  A^{\,p-k}\mid\mu_G,\qquad A=\mu_{K_{1,m}}=x^{m-1}(x^2-m),
\]
which is the two-hub identity's $A^{p-2}$ for $k=2$. At $k=3$, $m=3$, $p=5$ the graph has $23$
vertices, minimum degree three and vertex connectivity \emph{three}, and $\sqrt3$ is a root of
$\mu_G$ lying outside $\operatorname{spec}(T_G)$ (\texttt{code/c2\_attack.py}).
\end{proposition}

So neither stated hypothesis survives, and they fail for one reason rather than two: both constrain
the ambient graph, while the engine is a property of the divisor. A star's leaves have degree one in
$H$ whatever their degree in $G$, and enlarging the separator leaves $A$ untouched while raising
$\kappa$.

We do not push the reading further, because the measurements do not support it. Raising the
separator size $k$ raises the divisibility to $A^{p-k}$, at a cost of $p\ge k+1$ branches, but it
also raises every branch vertex's degree, and the two effects run against each other. Our scans
resolve neither, so we make no claim about $\kappa\ge4$ in either direction. What is established is
that the two hypotheses actually proposed as repairs are false.

The implication D3~$\Rightarrow$~C2 (\texttt{MinimumDegreeThreshold.C2\_of\_D3}) and the sharpness
statement (\texttt{threshold\_is\_three}) remain valid as implications, but both are now vacuous,
their hypotheses being false.

\subsection{What this leaves}\label{sec:leaves}

Conjecture~\ref{conj:gap} is false, and so are both hypotheses proposed to repair it. The three
refutations are one refutation: by Lemma~\ref{lem:divis} the branch polynomial $A$ divides $\mu_G$
whenever a separator of size $k$ is met by $p\ge k+1$ identical branches, and $A$ depends on the
branch alone. Minimum degree and vertex connectivity are conditions on $G$. Neither can see $A$,
and a star divisor is exactly a branch whose own degrees stay at one however large the degrees of
$G$ become. Any repair of Conjecture~\ref{conj:gap} must constrain the divisor.

We do not know where that leaves the biregular case, which is the setting the phenomenon actually
inhabits and is treated in \cite{PartII}. There the cover is an $(a,b)$-biregular tree, the gap is
the single interval $(0,|s-t|)$, and no construction of the present kind has reached it.

Two things are open and we state them as open rather than as programme. Whether some hypothesis on
the \emph{divisor} rescues the bound is untouched here; the natural candidate compares the
divisor's Heilmann--Lieb radius $2\sqrt{\Delta_H-1}$ with the band edge it would have to clear.
And whether higher connectivity eventually defeats the construction is not settled by our
measurements: raising the separator size $k$ raises every branch vertex's degree as well, and the
two effects run against each other. We looked, the picture did not resolve, and we make no claim
in either direction.
\section{Related work and formalization}

The $d$-matching polynomial and the interval $[-\rho,\rho]$ are due to Hall, Puder and Sawin; the
counterexample refuting the inclusion is due to Chris Hall \cite{HallCE} and is reproduced here
with his permission, with our own independent verification of its certificate. The
Angel--Friedman--Hoory ratio criterion \cite{AFH} is what certifies the gap.

The Lean~4/Mathlib sources are in \texttt{RamaLean/} of the accompanying repository, and carry no
\texttt{sorry} and no custom axioms. Formalized here: the closed form at first Betti number one
and its corollaries, the counterexample's algebra, the branch and two-cut factorizations, the
separation mechanism, and the minimum-degree threshold, including that minimum degree two does
not suffice and that three is therefore exact. Statements are labelled \texttt{VERIFIED},
\texttt{HEURISTIC} or \texttt{CONJECTURE} in the sources, and the labels are used with their
plain meaning: a \texttt{HEURISTIC} claim is supported by computation and is not proved.

\begin{thebibliography}{9}
\bibitem{GG} C.~Godsil, I.~Gutman, \emph{On the matching polynomial of a
graph}, in: Algebraic Methods in Graph Theory, 1981.
\bibitem{HPS} C.~Hall, D.~Puder, W.~Sawin, \emph{Ramanujan coverings of
graphs}, Adv.\ Math.\ \textbf{323} (2018), 367--410.
\bibitem{CGHRS} G.~Cochran, C.~Groothuis, A.~Herring, R.~Rohatgi,
E.~Stucky, \emph{A new (combinatorial) proof of the commutativity of
matching polynomials for cycles}, arXiv:1810.05889 (2018). Preprint; no journal
reference or DOI as of August 2026.
\bibitem{WWM} J.~Wan, W.~Wang, A.~Mohammadian, \emph{On the matching
polynomials of graphs with small number of cycles}, arXiv:2103.10580 (2021).
\bibitem{SFM} Y.-M.~Song, Y.-Z.~Fan, Z.~Miao, \emph{Hypergraph coverings and
Ramanujan hypergraphs}, arXiv:2310.01771 (2023).
\bibitem{HL} O.~J.~Heilmann, E.~H.~Lieb, \emph{Theory of monomer-dimer
systems}, Comm.\ Math.\ Phys.\ \textbf{25} (1972) 190--232.
\bibitem{God81} C.~D. Godsil, \emph{Matchings and walks in graphs},
J. Graph Theory \textbf{5} (1981), 285--297.

\bibitem{God93} C.~D. Godsil, \emph{Algebraic Combinatorics}, Chapman and Hall,
New York, 1993, Theorem 6.1.1 and \S5.6.
\bibitem{MS} H.~Mizuno, I.~Sato, \emph{Characteristic polynomials of some graph
coverings}, Discrete Math.\ \textbf{142} (1995) 295--298.
\bibitem{DFMRS} C.~Dalf\'o, M.~A.~Fiol, M.~Miller, J.~Ryan, J.~\v{S}ir\'a\v{n},
\emph{An algebraic approach to lifts of digraphs}, Discrete Appl.\ Math.\
\textbf{269} (2019) 68--76.
\bibitem{MSS} A.~Marcus, D.~Spielman, N.~Srivastava, \emph{Interlacing
families I: Bipartite Ramanujan graphs of all degrees}, Ann.\ of Math.\
\textbf{182} (2015), 307--325.
\bibitem{HallCE} C.~Hall, \emph{A simple bipartite counterexample to Conjecture 10:
  spectral localization of the ordinary matching polynomial}, verification note, personal
  communication, August 2026.
\bibitem{AFH} O.~Angel, J.~Friedman, S.~Hoory, \emph{The non-backtracking spectrum of the
  universal cover of a graph}, Trans.\ Amer.\ Math.\ Soc.\ \textbf{367} (2015), 4287--4318;
  arXiv:0712.0192.
\bibitem{LMST} W.~Li, M.~Magee, M.~Sabri, D.~Thomas, \emph{Eigenvalues of maximal
  abelian covers}, arXiv:2508.17332 (2025).
\bibitem{Sp} T.~J.~Spier, \emph{Eigenvalues of universal covers and the matching
  polynomial}, arXiv:2510.05041 (2025).
\bibitem{PartII} V.~Bonfioli, \emph{The biregular case of a spectral-inclusion
  conjecture: the weighted plane class and its obstructions}, companion paper, 2026.
\bibitem{Supp} V.~Bonfioli, \emph{Roots of $d$-matching polynomials: supplementary
  material}, companion note, 2026.
\bibitem{Xu} Z.~Xu, \emph{Note on mixed characteristic polynomials}, personal communication,
  August 2026.
\bibitem{MO} S.~Mohanty, R.~O'Donnell, \emph{X-Ramanujan graphs}, arXiv:1904.03500v2 (2019).
\end{thebibliography}

\end{document}
